\section{Supplementary Methods}
\subsection{Study design and clinical records}
This analysis used longitudinal records from a Chinese observational treat-to-target project. The supplied protocol specified clinician-directed systemic treatment rather than random allocation. We retained the initial treatment group for longitudinal comparisons and separately extracted subsequent dose, interval, stopping, switching and restart information. The analysed records came from CAMS and Huashan. The protocol's planned centre list and planned sample size were not substituted for the observed cohort.

The clinical dose-comparison population included patients aged at least 12 years with a known initial regimen. Records with unresolved simultaneous initial-dose assignments were excluded from dose-specific comparisons but retained in the source audit. Adult DLQI analyses used age at least 18. Repeated records were linked by the study patient identifier and nominal visit; a patient was never treated as a new independent person at each visit. The requested reporting visits were W0, W2, W4, W12 and W24 for PP-NRS; W0, W4, W12 and W24 for EASI, IGA, DLQI, BSA, SCORAD and eosinophils; and W0, W12 and W24 for IgE. Low-coverage intermediate visits were retained in the underlying data but were not inserted into these trajectories.

EASI was the recorded clinical reference outcome. Source discrepancies between regional sums and total EASI were audited rather than silently repaired. IGA grades were retained as recorded, including grade 5. Confirmation of the exact IGA version and PP-NRS recall window remains necessary. Nominal clinical/serum visit linkage does not by itself verify that the blood draw and examination occurred on the same calendar day; exact dates and a permissible matching window require source confirmation. Among the 152/89/75 model visits at W0/W12/W24, 8/2/0 had complete regional scores whose sum differed from total EASI by more than 0.15; the largest absolute gaps were 18.7/8.5/0.0 points, and 1/1/4 visits lacked a complete regional reference. Restricting evaluation to all four regions observed with total--sum difference at most 0.01 retained 140/86/71 visits and gave MAEs 3.17/1.32/1.11. This sensitivity evaluates unchanged saved predictions; it neither corrects the clinical records nor refits after training exclusions. Details are in \texttt{model\_reference\_QC\_audit.csv}.

\subsection{Clinical targets and treatment comparisons}
The week-2 target was PP-NRS $\leq4$. At week 4, the physician domain was EASI50 or an IGA decrease of at least two points, and the patient domain was PP-NRS $\leq4$ or adult DLQI $\leq10$. At week 12, the physician domain was EASI75 or IGA 0/1, and the patient domain was PP-NRS 0/1 or adult DLQI $\leq4$. At week 24, EASI $\leq5$ was an additional physician-domain criterion. The combined target required both domains. Because instrument identity and recall windows remain unresolved, these are exploratory targets based on the recorded scores. These are the investigators' operational definitions; they were not retrospectively relabelled as the newer IEC low-disease-activity or remission definitions.\citep{merola2026iec}

For an OR criterion, one observed positive component established success; a negative result required all components to be observed and negative. For an AND criterion, one observed negative component established failure; success required both positive components. Other cases remained unknown. Adult DLQI was an unknown component for ages 12--17, not a negative component and not removed from the OR rule. Thus adolescent itch success establishes the patient target; itch failure alone does not establish patient-target failure. The audit exhaustively verified the nine two-component truth states and reconciled every saved composite. Among 15 adolescents, the patient domain had 6/0/9 successes/failures/unknowns at W4, 4/0/11 at W12 and 4/0/11 at W24; the combined domain had 5/1/9, 4/5/6 and 4/3/8. An explicitly different PP-NRS-only adolescent rule would make five W12 and four W24 patient-domain unknowns determinate, including one W24 combined outcome. This is a sensitivity count, not a redefinition of the primary composite. Age-stratified counts are in \texttt{composite\_age\_evaluability\_audit.csv}. EASI response required an observed positive baseline. Individual thresholds were also analysed separately. Proportions used endpoint-specific observed denominators and Wilson 95\% confidence intervals. Overall three-group tests used conditional Monte Carlo Pearson tests with 49,999 draws, and pairwise comparisons used Fisher's exact test and risk-difference intervals. The three pairwise comparisons within an endpoint and visit formed a Holm family; exploratory FDR results were also retained. Overall tests were not included as an extra member of that pairwise family.

Baseline-adjusted binary target rates and risk differences used generalized overlap weights from a multinomial initial-treatment model with baseline EASI, age, sex, BMI and centre. We used 1,000 patient bootstrap resamples, stratified by initial treatment and refitting the propensity model, for percentile intervals. Two-sided P values used normal Wald tests based on bootstrap standard errors; three-group tests used a two-degree-of-freedom Wald statistic. The three individual pairwise comparisons formed a Holm family. Weighted results supplement the observed numerator and denominator. A revision W2 sensitivity added baseline PP-NRS to the propensity model. One patient lacked baseline PP-NRS; both the original-covariate model and the added-itch model were fitted on the remaining 162 patients, with the original 163-patient analysis retained. The same standardized logistic model (ridge parameter $C=1$), 1,000 stratified patient resamples and seed 20261018 were used for this sensitivity. The added-itch 100 mg--dupilumab contrast was 41.4 percentage points (95\% CI, 13.0--66.8; three-contrast Holm $P=0.0104$); the 100-minus-200 mg contrast was 6.7 points ($-11.3$ to 26.2). The largest endpoint-specific absolute standardized difference was 0.53 even after adding itch, so residual imbalance remained. Baseline weights were retained when calculating available-case balance and ESS; no claim of balanced follow-up selection is made. The table below gives unweighted worst/best bounds obtained by assigning every unknown outcome failure/success, alongside the available-case overlap-weight ESS. Bounds are sensitivity limits, not confidence intervals.
\begin{center}\small\textbf{Unweighted observed outcomes and initial-cohort bounds}\par\smallskip\begin{tabular}{llrrrr}\toprule Target & Initial group & Success/evaluable & Initial N & \shortstack{Observed\\rate, \%} & \shortstack{Unweighted missing-\\outcome bounds, \%} \\ \midrule
W2 itch & dupi & 15/27 & 65 & 55.6 & 23.1--81.5 \\
W2 itch & abro 100 mg & 34/37 & 56 & 91.9 & 60.7--94.6 \\
W2 itch & abro 200 mg & 33/36 & 42 & 91.7 & 78.6--92.9 \\
W24 combined & dupi & 26/36 & 65 & 72.2 & 40.0--84.6 \\
W24 combined & abro 100 mg & 22/24 & 56 & 91.7 & 39.3--96.4 \\
W24 combined & abro 200 mg & 20/23 & 42 & 87.0 & 47.6--92.9 \\
\bottomrule\end{tabular}\par\medskip\textbf{Matching baseline-overlap-weighted analysis}\par\smallskip\begin{tabular}{llrrr}\toprule Target & Initial group & \shortstack{Weighted available-\\case rate, \%} & \shortstack{Available-case\\weight ESS} & \shortstack{Weighted missing-\\outcome bounds, \%} \\ \midrule
W2 itch & dupi & 51.9 & 17.6 & 19.4--82.0 \\
W2 itch & abro 100 mg & 93.6 & 25.1 & 76.0--94.8 \\
W2 itch & abro 200 mg & 85.8 & 21.1 & 63.3--89.5 \\
W24 combined & dupi & 46.2 & 20.8 & 22.0--74.3 \\
W24 combined & abro 100 mg & 85.4 & 14.7 & 26.2--95.5 \\
W24 combined & abro 200 mg & 87.0 & 12.5 & 38.3--94.3 \\
\bottomrule\end{tabular}
\par\smallskip\parbox{.97\linewidth}{\footnotesize Bounds assign all unknown outcomes failure or success in the initial cohort. Weighted bounds retain the same baseline overlap weights in observed and unknown records. Weighted available-case rates use only endpoint-observed patients. ESS is the effective sample size of those weighted available cases, not a patient count or an incidence denominator. Bounds are sensitivity limits, not confidence intervals.}
\end{center}
 For generalized propensity scores $e_t(x)$, the weight is $h(x)/e_t(x)$, with $h(x)=1/\sum_t[1/e_t(x)]$. This targets a common overlap population. Covariate balance, propensity distributions and analysis-specific effective sample sizes are reported; weighting does not establish elimination of confounding.\citep{lili2019}

Continuous treatment comparisons used marginal regression with baseline outcome, treatment-by-visit terms, age, sex, BMI and centre. Generalized overlap weights addressed measured baseline differences in initial treatment selection. Patient-cluster robust covariance accounted for repeated outcomes, and a singular-value decomposition check prevented reporting non-estimable treatment-by-visit cells. IgE and eosinophils used log$(1+x)$; back-transformed estimates are not arithmetic means. Unweighted adjustment, group balance, observation patterns and supplementary dupilumab-versus-pooled-abrocitinib contrasts were retained. Pooling used the observed dose mixture and did not erase the two initial abrocitinib groups. The comparisons estimate adjusted observational associations, not randomized fixed-dose effects. Nominal eligibility based on the registry filename date placed all initial patients beyond W24, but this administrative sensitivity is not an approved database lock or verified visit window. Reasons for absent assessments were unavailable.

Treatment reductions were assigned from recorded regimen intervals, with assessment/decision ordering marked uncertain at shared visit boundaries. A later reduction was not moved to an earlier visit. W12 response and subsequent regimen changes defined descriptive maintenance pathways. Analyses used observed W24 EASI75, explicit missing-outcome bounds and fixed-patient trajectories where all three visits were observed. This approach distinguishes treatment initiation from the selected maintenance pathways, consistent with the need to specify how intercurrent treatment changes enter the question being estimated.\citep{iche9r1} It does not identify a counterfactual outcome under a different regimen. At W0, two source records labelled initial 100 mg explicitly listed a 200 mg starting regimen, and one labelled initial 200 mg listed 100 mg. These are conflicts between the source group and regimen fields, not demonstrated post-baseline switches. The group labels were retained, and \texttt{exposure\_W0\_source\_adjudication.csv} preserves their counts separately from unknown exposure. Clinical-record adjudication is still required.

\subsection{Serum measurements and molecular analyses}
The laboratory source reports identify serum samples analysed with Luminex and MSD platforms. Luminex used an R\&D Systems LXSAHM-13 panel and Luminex-200; the MSD reports identify K15067M-1/K15067M-2 kits and an SQ120 instrument. Reports describe 50-$\mu$L serum and twofold dilution. The service-batch sample count was not used as the study sample size. Assay lots, run identifiers and source links remain in the restricted provenance records.

The analysis used the numeric integration export rather than the vendor's filtered display matrix. Positive export numerics recorded in pg/mL were log$_2$ transformed once, unchanged; the analysis applied no additional twofold dilution correction. This reproduces the fitted scale rather than verifying how the laboratory corrected its export for dilution. Explicit censoring strings and nonpositive values were not converted to measured positive concentrations. The vendor's filtered output can replace values below range with zero and values above range with an upper bound; those conventions must not be interpreted as exact measured concentrations. A sensitivity analysis represented IL-5 values below twice the software calibration lower bound at that boundary and repeated the model workflow. The software boundary is not an independently validated clinical quantification limit. Cross-platform bridging, technical-repeat precision and lot effects require additional assay documentation.

All 23 measured markers were retained. Molecular changes were formed within each patient and visit pair before group summaries, avoiding subtraction of averages from different patient sets. We analysed W0--W12, W12--W24 and W0--W24 separately. Mean log$_2$ changes were accompanied by t-based confidence intervals when at least five pairs were available. Smaller sets were descriptive. Two-sided paired mean-change sign-flip tests retained zero changes and ties. They enumerated all $2^n$ signs when $n\leq15$; otherwise they used 19,999 random signs (seed 20261002), counting exceedances at the observed absolute mean with a plus-one numerator and denominator. All-zero changes gave $P=1$. The sign-flip null assumes independent patients and sign-exchangeable paired changes (symmetry about zero); it is not a randomized-treatment test. Benjamini--Hochberg adjustment covered the full 23-marker, three-treatment, three-window family: 207 planned cells, of which 193 had estimable P values. Non-estimable cells remained missing and were not counted as tests. Between-treatment change models adjusted for starting concentration, baseline EASI, age, sex, BMI and centre, with HC3 covariance. Centre-specific and assay-batch sensitivity analyses were retained. Each between-treatment molecular correction family included all marker, window and pairwise contrasts within one centre scope and analysis variant; scopes and variants were not pooled. Molecular--clinical change correlations formed their own panel-wide family across markers, clinical outcomes, treatments and windows, and response-depth contrasts formed a separate family. The global molecular treatment-by-visit family had 23 planned marker tests, with 21 estimable. A source-table inventory records the defined families and finite-test counts. Spearman correlations related molecular changes to clinical changes within treatment groups; baseline regional correlations were assessed separately. Sample sizes, intervals and multiplicity families accompany each source table.

Regional analyses used source-consistent regional totals. Regional EASI contributions were divided by their theoretical maxima (7.2, 14.4, 21.6 and 28.8 for head/neck, upper limbs, trunk and lower limbs). Percentage improvement required a positive regional starting value; zero-start exclusions were counted. Because very low W12 values can produce extreme W12--W24 percentage changes, the corresponding figure instead shows change as a percentage of the theoretical regional maximum. Both definitions are retained in the tables. Regional treatment contrasts used unadjusted Welch comparisons and HC3 models adjusted for starting regional severity, baseline total EASI, age, sex, BMI and centre, with FDR over the 48 region-by-window-by-contrast tests in each analysis/outcome family. Whole-body PP-NRS change was not labelled as regional itch.

\subsection{Common EASI estimator and validation}
Eligibility required age at least 12, observed EASI and BSA, and observed positive CCL17, CCL26 and IL-5 concentrations. Ambiguous initial-dose labels did not exclude otherwise eligible model records, because treatment was not a required input. The three-marker panel arose from earlier exploratory work and was fixed for this revision. No clinical outcome was imputed. Model development used W0 and W12; W24 did not enter fitting or selection. Because earlier project analyses had already examined W24, its results are described as within-cohort temporal application rather than untouched external validation.

All observations from one patient remained in the same outer fold. Five-fold validation was repeated twice using the established patient partitions. Each training fold contained development visits only. Three-fold patient-level inner validation selected the ridge penalty from 0.1, 1, 10 and 100, minimizing mean per-patient absolute error. Each patient received equal total fitting weight. Feature centring and scaling used training data only. Three common structures were considered: an additive BSA-plus-three-marker model; the same model with quadratic BSA; and BSA multiplied by a linear molecular component without a free intercept. Structure and penalty selection were nested together for the primary performance estimate. No treatment, phase, time, previous score, age or BMI was required at deployment.

After full-development inner selection, the fitted common equation was
\begin{align*}
\widehat{\mathrm{EASI}}=\operatorname{clip}_{[0,72]}\{\mathrm{BSA}\,[&0.512721579
 +0.009847080\log_2(\mathrm{CCL17})\\
 &+0.009646929\log_2(\mathrm{CCL26})
 -0.008270558\log_2(\mathrm{IL\mbox{-}5})]\}.
\end{align*}
BSA is entered as a percentage from 0 to 100, and concentrations must use the development export's recorded pg/mL numeric scale, with log$_2$ applied once. The implementation rejects missing, nonfinite or nonpositive molecular inputs and invalid BSA values. The area constraint returns zero when BSA is zero. It is a statistical constraint reflecting absent observed lesion extent, not proof of absent molecular inflammation. The released parameter file retains full precision. The main-text display multiplies the molecular coefficients by 100 and uses BSA/100, allowing two decimal places without dropping the smaller molecular terms. An explicit rounding audit gave nested-selection W0/W12/W24 MAEs of 3.194/1.298/1.193, compared with 3.194/1.297/1.192 at full precision. All primary results use full precision; the rounded equation is a reading aid. For the final fitted equation, maximum absolute arithmetic differences were 0.075550, 0.024807 and 0.018474 EASI points at W0/W12/W24; these are rounding differences, not clinical prediction errors. Response classifications at EASI50/75/90 were unchanged in the audit, but one of 150 week-24 fold-level evaluations crossed the absolute EASI 7 threshold, so threshold-based use should retain full precision. This single final fit has not received independent validation; primary metrics evaluate the nested development procedure that selected it.

Comparators used the same patients and folds: BSA alone, the three markers alone, the additive common model, and the earlier stage-interaction model. We evaluated MAE, RMSE, signed error, Spearman correlation, Lin's concordance correlation, $R^2$, calibration intercept and slope, severity-band agreement and proportions within one and two EASI points. Metrics were calculated within each repetition before averaging, not from averaged predictions. Patient bootstrap intervals used 2,000 resamples of held-out errors. The same sampled patient indices were used in every repetition and comparator; metrics were scored within each repetition before averaging. These intervals describe uncertainty conditional on the saved cross-validation predictions: neither nested selection nor fitting was repeated, and earlier adaptive marker exploration is not covered. Calibration intervals resampled patients in the regression of observed on predicted EASI. The primary nested-procedure minus area-only MAE differences were 0.12 (95\% CI, 0.04--0.20) at W12 and 0.17 (0.11--0.24) at W24. These paired estimates differ from the earlier fixed-structure comparison and do not establish molecular added accuracy. All comparator and fixed-structure results remain in the source tables. Descriptive agreement limits were mean error plus or minus 1.96 error standard deviations.\citep{bland2007}

Severity-band agreement used four categories: EASI $\leq7$, $>7$ to $\leq21$, $>21$ to $\leq50$ and $>50$. Error subgroup analyses instead combined the two highest categories, using three strata: EASI $\leq7$, $>7$ to $\leq21$ and $>21$. Binary sensitivity and specificity used EASI $>7$ and are not interchangeable with every original p-EASI classification definition. Leave-centre-out and leave-initial-treatment-group-out analyses removed all records from held-out patients. Changes were paired within patient, repetition and fold before averaging. Each of W0--W12, W0--W24 and W12--W24 was evaluated separately; the windows share patients. EASI50/75/90 classification required positive observed and estimated starting scores, with exclusions counted. Observed response used observed scores at both visits, and estimated response used estimated scores at both visits. Scatter plots and any integer confusion matrices use the first outer repetition; the primary response forest and principal metrics average repetition-specific values. They must not be interpreted as one common confusion matrix. EASI75 sensitivity used 49 and 54 observed responders at W12 and W24; specificity used 36 and 13 observed nonresponders, respectively, in both repetitions. Patient-bootstrap intervals for sensitivity, specificity and overall agreement are provided in \texttt{primary\_EASI75\_uncertainty.csv}. The empirical cumulative absolute-error curve is the average of the two repetition-specific proportions. Pointwise bootstrap intervals carry all saved predictions from each sampled patient together and do not refit the development procedure. Applying the published formulas used the published rounded coefficients without local refitting, with local IL2RA as a nominal sIL-2R proxy and identical eligible patients for both methods. Original p-EASI outputs were initially left unbounded to preserve the published arithmetic, whereas local estimates were clipped to 0--72; matched clipping of the original outputs was examined separately. Patient-paired bootstrap intervals were calculated for differences between the current and published-formula metrics, retaining both repetitions for each sampled patient. These intervals characterize the local comparison; they do not remove the difference in inputs or establish assay equivalence.\citep{thijs2017,thijs2019,bakker2020}

\subsection{Independent auxiliary target models}
PP-NRS and adult DLQI were fitted separately with original-scale or log$(1+x)$ ridge regression, using patient-level inner validation to select the link and penalty. Predictions were bounded to 0--10 and 0--30, respectively. Recorded IGA used a penalized proportional-odds model over the six observed grades, with ordered thresholds and six category probabilities. Normalized ranked probability score was the primary ordinal loss. All three targets used the same four-input panel for a practical shared measurement set, with BSA-only, molecular-only and training-only constant comparators. Neither EASI nor the target's current observed score was an input. Development and temporal application followed the patient-level W0/W12 and W24 separation used above. This limited shared-panel study does not claim to select the optimal molecular set independently for each target. Details and full estimates are in the supplement.

\subsection{Reproducibility and reporting}
The model analysis used Python 3.9.6, NumPy 2.0.2, pandas 2.3.3, SciPy 1.13.1 and scikit-learn 1.6.1. Exact dependency versions, analysis specifications, input/output hashes and execution logs accompany the code. Clinical and prediction reporting follow the relevant STROBE and TRIPOD+AI domains; applicability and risk-of-bias issues are recorded explicitly.\citep{vonelm2007strobe,collins2024tripod,moons2025probast} No non-significant difference, high correlation or overlapping mean trajectory was taken as proof of individual interchangeability.
